\chapter{ЭКСПЕРИМЕНТАЛЬНАЯ ОЦЕНКА И АНАЛИЗ РЕЗУЛЬТАТОВ}

\section{Методика проверки}

Экспериментальная оценка выполнялась после обновления модели, используемой веб-интерфейсом. Поэтому в качестве основной модели рассматривалась не ранняя дипломная модель, а актуальная модель \texttt{v2\_medium\_rtx2060\_fast}, которая автоматически выбирается при запуске веб-приложения:

\begin{itemize}
\item checkpoint: \texttt{results/v2\_medium\_rtx2060\_fast/checkpoints/structured\_pipeline.pt};
\item конфигурация: \texttt{configs/pipeline\_v2\_medium\_rtx2060\_fast.json};
\item режим генерации: фрагментированная генерация по 4 такта на FASTA-фрагмент;
\item seed для воспроизводимости: 42.
\end{itemize}

Для сравнения использовались две дополнительные системы: предыдущая дипломная модель \texttt{thesis\_final\_run\_regularized} и случайный baseline. Оценка проводилась на 12 FASTA-фрагментах длиной 1800 bp, полученных из файла \texttt{data/fasta/training/thesis\_reference\_genomes.fa}. Такой набор соответствует прежней дипломной оценке, но позволяет проверить именно тот режим генерации, который доступен пользователю через веб-интерфейс.

\section{Метрики оценки}

\subsection{Технические метрики}

\textbf{Valid MIDI rate} — доля корректно сгенерированных MIDI-файлов.

\textbf{Melody note count} — количество нот в мелодии. При сравнении моделей разной длины эта метрика используется только как вспомогательная.

\textbf{Note density per bar} — плотность нот на такт. Это нормированная метрика, поэтому она важнее абсолютного количества нот при сравнении 4-тактовой и 8-тактовой генерации.

\textbf{Pitch range} — диапазон высот в полутонах.

\textbf{Unique pitches} — количество уникальных высот.

\textbf{Pitch class entropy} — энтропия распределения pitch classes.

\textbf{Chord change rate} — частота смены аккордов.

\textbf{Chord-tone ratio} — доля мелодических нот, принадлежащих текущему аккорду. Это ключевая метрика для оценки согласованности мелодии и гармонии. Для фрагментированной генерации значение вычислялось по фактическим MIDI-аккордам, так как metadata веб-режима не хранит полный список \texttt{generated\_harmony\_bars}.

\textbf{Self-similarity} — мера повторяемости локальных мелодических паттернов.

\subsection{Baseline}

Random baseline создаёт валидный двухдорожечный MIDI-файл, но не использует биологические признаки и обученные зависимости:
\begin{itemize}
\item случайная гармония;
\item случайная мелодия;
\item фиксированный seed для воспроизводимости;
\item та же процедура расчёта метрик, что и для моделей.
\end{itemize}

\section{Результаты обучения}

\subsection{Динамика loss}

На рисунке~\ref{fig:training} показаны кривые обучения актуальной web-модели и предыдущей дипломной модели. Актуальная модель обучалась на большем наборе пар и на более коротких 4-тактовых сегментах, что упростило локальную задачу генерации.

\begin{figure}[htbp]
\centering
\includegraphics[width=0.95\textwidth]{training_loss_curves.png}
\caption{Динамика train и validation loss для актуальной и предыдущей моделей}
\label{fig:training}
\end{figure}

\subsection{Сравнение loss}

\begin{table}[htbp]
\centering
\caption{Сравнение loss актуальной web-модели и предыдущей дипломной модели}
\label{tab:model-loss-comparison}
\begin{tabular}{|l|c|c|c|}
\hline
\textbf{Метрика} & \textbf{Web-модель} & \textbf{Предыдущая модель} & \textbf{Улучшение} \\
\hline
Harmony best val loss & 0.1454 & 0.2112 & 31.1\% \\
\hline
Melody best val loss & 0.1566 & 0.2556 & 38.8\% \\
\hline
Harmony test loss & 0.1389 & 0.2187 & 36.5\% \\
\hline
Melody test loss & 0.1650 & 0.2561 & 35.6\% \\
\hline
\end{tabular}
\end{table}

Снижение validation и test loss показывает, что актуальная модель лучше воспроизводит распределение обучающих последовательностей. Это не является прямой оценкой художественного качества, но является важным техническим индикатором качества генеративной модели.

\section{Результаты генерации}

\subsection{Сравнение моделей и baseline}

В таблице~\ref{tab:evaluation} приведено сравнение актуальной web-модели, предыдущей дипломной модели и random baseline на одних и тех же 12 FASTA-фрагментах.

\begin{table}[htbp]
\centering
\caption{Сравнение метрик генерации на 12 FASTA-фрагментах}
\label{tab:evaluation}
\begin{adjustbox}{max width=\textwidth}
\begin{tabular}{|l|c|c|c|}
\hline
\textbf{Метрика} & \textbf{Web-модель} & \textbf{Предыдущая модель} & \textbf{Random baseline} \\
\hline
Valid MIDI rate & 12/12 & 12/12 & 12/12 \\
\hline
Melody notes (mean) & 20.83 & 44.25 & 42.33 \\
\hline
Note density per bar & 5.21 & 4.92 & 5.29 \\
\hline
Pitch range & 14.25 & 15.92 & 34.17 \\
\hline
Unique pitches & 8.33 & 11.92 & 22.83 \\
\hline
Pitch class entropy & 2.50 & 2.91 & 3.28 \\
\hline
Chord change rate & 0.917 & 0.781 & 0.691 \\
\hline
\textbf{Chord-tone ratio} & \textbf{0.6710} & \textbf{0.6073} & \textbf{0.4043} \\
\hline
Self-similarity & 0.166 & 0.131 & 0.109 \\
\hline
\end{tabular}
\end{adjustbox}
\end{table}

На рисунке~\ref{fig:comparison} показано визуальное сравнение основных метрик. Для графика выбраны показатели, которые наиболее прямо описывают гармоническую согласованность, плотность и мелодический диапазон.

\begin{figure}[htbp]
\centering
\includegraphics[width=0.95\textwidth]{evaluation_metric_comparison.png}
\caption{Сравнение метрик актуальной web-модели, предыдущей модели и baseline}
\label{fig:comparison}
\end{figure}

\subsection{Интерпретация результатов}

\textbf{Chord-tone ratio.} Актуальная web-модель показала значение 0.6710 против 0.6073 у предыдущей дипломной модели и 0.4043 у random baseline. Относительно baseline прирост составляет 66.0\%, относительно предыдущей модели — 10.5\%. Это означает, что мелодия чаще попадает в состав текущего аккорда, то есть лучше согласуется с гармоническим контекстом.

\textbf{Loss.} По всем четырём loss-показателям актуальная web-модель лучше предыдущей: harmony validation loss ниже на 31.1\%, melody validation loss — на 38.8\%, harmony test loss — на 36.5\%, melody test loss — на 35.6\%. Это связано с большим числом обучающих пар и с тем, что 4-тактовые сегменты задают более простую локальную задачу.

\textbf{Длина и плотность.} Абсолютное число нот у web-модели ниже, потому что она генерирует 4-тактовые фрагменты, а предыдущая модель — более длинные 8-тактовые фразы. Поэтому корректнее сравнивать \texttt{note\_density\_per\_bar}: 5.21 у web-модели против 4.92 у предыдущей модели.

\textbf{Pitch range и unique pitches.} Web-модель использует более сфокусированный мелодический словарь: pitch range 14.25 против 34.17 у baseline, unique pitches 8.33 против 22.83. Это не следует трактовать как универсальное улучшение по всем музыкальным критериям. Скорее это показывает, что модель удерживает мелодию в более компактном регистровом и ладовом пространстве.

\section{Анализ сгенерированных композиций}

На рисунке~\ref{fig:profile} представлен профиль 12 MIDI-фрагментов, сгенерированных актуальной web-моделью. Видно, что количество нот меняется от фрагмента к фрагменту, но при этом chord-tone ratio остаётся выше случайного baseline.

\begin{figure}[htbp]
\centering
\includegraphics[width=0.95\textwidth]{generated_midi_profile.png}
\caption{Профиль MIDI-фрагментов, сгенерированных актуальной web-моделью}
\label{fig:profile}
\end{figure}

Каждый MIDI-фрагмент содержит:
\begin{itemize}
\item две дорожки: harmony и melody;
\item 4 такта музыки;
\item монофоническую мелодию;
\item аккордовую гармонию;
\item metadata с checkpoint, конфигурацией, длиной входного фрагмента и параметрами генерации.
\end{itemize}

\section{Ограничения разработанного решения}

\subsection{Технические ограничения}

\begin{itemize}
\item Модель ограничена символическим MIDI-представлением и не генерирует аудио напрямую.
\item Ограничение GPU памяти (6GB) не позволило использовать крупные ESM embeddings в финальном workflow.
\item ViennaRNA отключена из-за медленного preprocessing.
\item Веб-режим использует фрагментированную генерацию: длинная последовательность превращается в набор 4-тактовых фраз, а не в единую длинную форму.
\end{itemize}

\subsection{Методологические ограничения}

\begin{itemize}
\item Система не доказывает причинную связь между биологией и музыкой.
\item Музыка зависит от стиля обучающего корпуса POP909.
\item Биологические признаки используются как управляющий сигнал, а не как прямое отображение нуклеотидов в ноты.
\item Метрики оценивают технические свойства MIDI, но не заменяют экспертное или пользовательское прослушивание.
\end{itemize}

\subsection{Музыкальные ограничения}

\begin{itemize}
\item Мелодия строго монофоническая.
\item Гармония ограничена фиксированным набором качеств аккордов.
\item Динамика и артикуляция представлены упрощённо.
\item Фрагменты требуют дальнейшей музыкальной аранжировки, если нужно получить полноценную композицию.
\end{itemize}

\section{Оценка достоверности результатов}

\subsection{Воспроизводимость}

Все результаты воспроизводимы:
\begin{itemize}
\item фиксированный seed 42;
\item сохранённые checkpoint;
\item документированные конфиги;
\item отчёты \texttt{model\_comparison.json} и \texttt{model\_comparison.md};
\item MIDI-файлы для web-модели, предыдущей модели и baseline.
\end{itemize}

\subsection{Валидность метрик}

Метрики подтверждают:
\begin{itemize}
\item корректность генерации: valid MIDI rate 100\%;
\item согласованность мелодии и гармонии: chord-tone ratio выше baseline и предыдущей модели;
\item контролируемый мелодический диапазон;
\item воспроизводимое сравнение нескольких моделей на одном наборе FASTA-фрагментов.
\end{itemize}

\subsection{Ограничения интерпретации}

Результаты \textbf{не доказывают}, что:
\begin{itemize}
\item в ДНК содержится музыка;
\item существует причинная связь между генами и музыкальными структурами;
\item сгенерированная музыка имеет биологический смысл;
\item технические метрики полностью описывают художественное качество.
\end{itemize}

Результаты \textbf{подтверждают}, что:
\begin{itemize}
\item биологические признаки могут использоваться как conditioning signal для генератора музыки;
\item иерархическая архитектура создаёт валидные двухдорожечные MIDI-файлы;
\item актуальная web-модель улучшила гармоническую согласованность мелодии по сравнению с предыдущей моделью и random baseline;
\item фрагментированная генерация удерживает каждый участок в пределах обучающего распределения модели.
\end{itemize}

\section*{Выводы по главе 3}

В третьей главе представлена повторная экспериментальная оценка системы после обновления модели, используемой веб-интерфейсом. Актуальная модель \texttt{v2\_medium\_rtx2060\_fast} показала лучшие loss-показатели по сравнению с предыдущей дипломной моделью: harmony validation loss 0.1454 против 0.2112 и melody validation loss 0.1566 против 0.2556.

Сравнение генерации на 12 FASTA-фрагментах показало, что актуальная web-модель создаёт валидные MIDI-файлы и повышает chord-tone ratio до 0.6710. Это выше, чем у предыдущей модели (0.6073), и значительно выше random baseline (0.4043). При этом некоторые показатели, например unique pitches и pitch class entropy, следует трактовать как tradeoff: модель использует более сфокусированный мелодический словарь, а не превосходит альтернативы по каждому возможному дескриптору.

Таким образом, обновление модели улучшило техническую согласованность мелодии и гармонии, а также сделало оценку ближе к реальному пользовательскому режиму веб-приложения.
